\section{Sujet n\textsuperscript{o}\,29}
\noindent\begin{minipage}{0.5\linewidth}\footnotesize Office du Baccalauréat du Cameroun\end{minipage}%
\hfill\begin{minipage}{0.45\linewidth}\footnotesize\raggedleft Examen : \textbf{Baccalauréat} · Série : \textbf{C/E}\\
Épreuve : \textbf{Mathématiques} · Durée : \textbf{4\,h} · Coef. : \textbf{7}\end{minipage}
\par\smallskip{\color{onbo}\rule{\linewidth}{1pt}}

\subsection*{Partie A — Évaluation des ressources \hfill (15 points)}

\noindent\textbf{Exercice 1.} \hfill\textit{(3 points)}\\
Dans une imprimerie de Yaoundé, une presse numérique produit des affiches dont
$8\,\%$ présentent un défaut de calage des couleurs. Pour le contrôle qualité, on prélève
au hasard un échantillon de $n=6$ affiches dans une production très importante (le tirage
est assimilé à un tirage avec remise). On note $X$ la variable aléatoire égale au nombre
d'affiches défectueuses parmi les $6$ prélevées.
\begin{enumerate}[label=\arabic*)]
\item Justifier que $X$ suit une loi binomiale dont on précisera les paramètres. \hfill\textit{(0,75 pt)}
\item Calculer, à $10^{-4}$ près, la probabilité que l'échantillon ne contienne aucune
affiche défectueuse. \hfill\textit{(0,75 pt)}
\item Calculer, à $10^{-4}$ près, la probabilité que l'échantillon contienne exactement
une affiche défectueuse. \hfill\textit{(0,75 pt)}
\item Déterminer l'espérance mathématique $E(X)$ et la variance $V(X)$, puis interpréter
$E(X)$. \hfill\textit{(0,75 pt)}
\end{enumerate}

\smallskip
\noindent\textbf{Exercice 2.} \hfill\textit{(3 points)}\\
L'atelier comporte quatre postes de travail $A$ (composition), $B$ (impression),
$C$ (façonnage) et $D$ (expédition) reliés par des convoyeurs à sens unique :
de $A$ vers $B$, de $A$ vers $C$, de $B$ vers $C$, de $C$ vers $D$ et de $B$ vers $D$.
On note $M$ la matrice d'adjacence du graphe orienté associé, les sommets étant rangés
dans l'ordre $A,B,C,D$.
\begin{enumerate}[label=\arabic*)]
\item Écrire la matrice d'adjacence $M$. \hfill\textit{(1 pt)}
\item Calculer $M^2$. \hfill\textit{(1 pt)}
\item En déduire le nombre de chemins de longueur~$2$ allant de $A$ à $D$ et les citer. \hfill\textit{(1 pt)}
\end{enumerate}

\smallskip
\noindent\textbf{Exercice 3.} \hfill\textit{(4 points)}\\
On considère les fonctions $f$ et $g$ définies sur $\R$ par $f(x)=x^2$ et $g(x)=x+2$, de
courbes respectives $(\mathcal{C}_f)$ et $(\mathcal{C}_g)$ dans un repère orthonormé.
Elles modélisent deux profils de coût d'impression (en unités) en fonction du tirage~$x$.
\begin{enumerate}[label=\arabic*)]
\item Déterminer les abscisses des points d'intersection de $(\mathcal{C}_f)$ et
$(\mathcal{C}_g)$. \hfill\textit{(1 pt)}
\item Étudier le signe de $g(x)-f(x)$ sur l'intervalle $[-1\,;\,2]$. \hfill\textit{(1 pt)}
\item Calculer, en unités d'aire, l'aire~$\mathcal{A}$ du domaine du plan compris entre
les deux courbes sur $[-1\,;\,2]$. \hfill\textit{(1,5 pt)}
\item Donner la valeur exacte de $\mathcal{A}$, puis une valeur approchée à $10^{-2}$ près. \hfill\textit{(0,5 pt)}
\end{enumerate}

\smallskip
\noindent\textbf{Exercice 4.} \hfill\textit{(5 points)}\\
Le plan est muni d'un repère orthonormé direct $(O\,;\,\vec\imath,\vec\jmath)$ (unité :
le mètre). On modélise la disposition de trois machines par les points $A(0\,;\,0)$,
$B(6\,;\,0)$ et $C(2\,;\,4)$.
\begin{enumerate}[label=\arabic*)]
\item Calculer le produit scalaire $\vec{AB}\cdot\vec{AC}$, puis une valeur approchée au
degré près de l'angle $\widehat{BAC}$. \hfill\textit{(1,25 pt)}
\item On considère le barycentre $G$ du système $\{(A,1),(B,2),(C,1)\}$. Calculer les
coordonnées de $G$. \hfill\textit{(1,25 pt)}
\item Déterminer et construire l'ensemble $(\Gamma)$ des points $M$ du plan vérifiant
$\vec{MA}+2\vec{MB}+\vec{MC}=4\,\vec{MA}$. \hfill\textit{(1,25 pt)}
\item Déterminer l'ensemble $(\mathcal{E})$ des points $M$ tels que
$\big\|\vec{MA}+2\vec{MB}+\vec{MC}\big\|=8$. \hfill\textit{(1,25 pt)}
\end{enumerate}

\subsection*{Partie B — Évaluation des compétences \hfill (5 points)}

\noindent\textit{Madame Ngo Bell dirige une maison d'édition à Yaoundé et te sollicite, en
tant qu'élève de Terminale~C, pour trois décisions. Pour un tirage de $x$ centaines de
livres, le coût total de fabrication (en milliers de francs) est
$C(x)=x^2-8x+25$. Par ailleurs, elle a lancé un abonnement scolaire dont le nombre
d'abonnés, parti de $u_0=500$, augmente chaque année de \SI{15}{\percent} grâce au
bouche-à-oreille : $u_{n+1}=1{,}15\,u_n$. Enfin, à la reliure, \SI{4}{\percent} des livres
sont mal collés ; un libraire prélève $3$ livres au hasard dans un grand stock.}

\begin{enumerate}[label=\textbf{Tâche \arabic*.},leftmargin=2.4em]
\item Déterminer le tirage qui rend le coût total \emph{minimal} et la valeur de ce coût
minimal. \hfill\textit{(1,5 pt)}
\item Déterminer l'année à partir de laquelle le nombre d'abonnés dépassera $1000$. \hfill\textit{(1,5 pt)}
\item Déterminer la probabilité qu'\emph{au moins un} des $3$ livres prélevés soit mal
collé (arrondir à $10^{-3}$). \hfill\textit{(1,5 pt)}
\end{enumerate}
\par\smallskip{\footnotesize\itshape Présentation générale : 0,5 point.}

% ════════════════════════════════════════════════════
\corrige{du sujet 29}

\noindent\textbf{Partie A — Exercice 1.}
1) On répète $n=6$ fois, de façon indépendante (tirage avec remise), une épreuve de
Bernoulli « l'affiche est défectueuse » de paramètre $p=0{,}08$ ; $X$ compte les succès,
donc $X\sim\mathcal{B}(6\,;\,0{,}08)$.
2) $P(X{=}0)=(0{,}92)^6\approx0{,}6064$.
3) $P(X{=}1)=\binom61(0{,}08)(0{,}92)^5=6\times0{,}08\times0{,}92^5\approx6\times0{,}08\times0{,}6591\approx0{,}3164$.
4) $E(X)=np=6\times0{,}08=0{,}48$ ; $V(X)=np(1-p)=6\times0{,}08\times0{,}92=0{,}4416$.
En moyenne, sur de tels échantillons de $6$ affiches, on observe environ $0{,}48$ affiche
défectueuse.

\noindent\textbf{Exercice 2.}
1) Arcs : $A\to B$, $A\to C$, $B\to C$, $B\to D$, $C\to D$. D'où, dans l'ordre $A,B,C,D$,
$$M=\begin{pmatrix}0&1&1&0\\0&0&1&1\\0&0&0&1\\0&0&0&0\end{pmatrix}.$$
2) $M^2=\begin{pmatrix}0&0&1&2\\0&0&0&1\\0&0&0&0\\0&0&0&0\end{pmatrix}$.
(Par exemple ligne $A$ : $A\to B\to C$ donne le coefficient en colonne $C$ ; $A\to B\to D$
et $A\to C\to D$ donnent $2$ en colonne $D$.)
3) Le coefficient $(A,D)$ de $M^2$ vaut $2$ : il y a \textbf{deux} chemins de longueur~$2$
de $A$ à $D$, à savoir $A\to B\to D$ et $A\to C\to D$.

\noindent\textbf{Exercice 3.}
1) $f(x)=g(x) \iff x^2=x+2 \iff x^2-x-2=0 \iff (x-2)(x+1)=0$ : abscisses $x=-1$ et $x=2$.
2) $g(x)-f(x)=-(x^2-x-2)=-(x-2)(x+1)$. Sur $[-1\,;\,2]$, $(x-2)\le0$ et $(x+1)\ge0$, donc
$-(x-2)(x+1)\ge0$ : $g(x)-f(x)\ge0$ sur $[-1\,;\,2]$ (la courbe de $g$ est au-dessus de
celle de $f$).
3) L'aire vaut $\mathcal{A}=\displaystyle\int_{-1}^{2}\big(g(x)-f(x)\big)\dd x
=\int_{-1}^{2}(-x^2+x+2)\dd x
=\Big[-\dfrac{x^3}{3}+\dfrac{x^2}{2}+2x\Big]_{-1}^{2}$.
En $x=2$ : $-\tfrac83+2+4=\tfrac{10}{3}$. En $x=-1$ : $\tfrac13+\tfrac12-2=-\tfrac76$.
Donc $\mathcal{A}=\tfrac{10}{3}-\big(-\tfrac76\big)=\tfrac{10}{3}+\tfrac76=\tfrac{20}{6}+\tfrac{7}{6}=\tfrac{27}{6}=\tfrac92$.
4) $\mathcal{A}=\dfrac{9}{2}=4{,}5$ unités d'aire.

\noindent\textbf{Exercice 4.}
1) $\vec{AB}=(6\,;\,0)$, $\vec{AC}=(2\,;\,4)$, donc $\vec{AB}\cdot\vec{AC}=6\times2+0\times4=12$.
$\cos\widehat{BAC}=\dfrac{12}{\|\vec{AB}\|\,\|\vec{AC}\|}=\dfrac{12}{6\times\sqrt{20}}
=\dfrac{12}{6\times2\sqrt5}=\dfrac{1}{\sqrt5}\approx0{,}4472$ ; $\widehat{BAC}\approx63°$.
2) Masse totale $1+2+1=4\neq0$ : $G$ existe et
$x_G=\dfrac{1\cdot0+2\cdot6+1\cdot2}{4}=\dfrac{14}{4}=\dfrac72$,
$y_G=\dfrac{1\cdot0+2\cdot0+1\cdot4}{4}=\dfrac{4}{4}=1$. Donc $G\big(\tfrac72\,;\,1\big)$.
3) Pour tout $M$, $\vec{MA}+2\vec{MB}+\vec{MC}=4\,\vec{MG}$ (réduction barycentrique).
L'équation devient $4\vec{MG}=4\vec{MA}$, soit $\vec{MG}=\vec{MA}$, donc $\vec{MA}-\vec{MG}=\vec0$,
c'est-à-dire $\vec{GA}=\vec0$… ce qui est impossible puisque $G\neq A$ ; vérifions :
$4\vec{MG}=4\vec{MA}\iff \vec{MG}=\vec{MA}\iff \vec{AG}=\vec0$, faux. Donc $(\Gamma)=\varnothing$.
(L'équation n'a aucune solution car elle équivaut à $G=A$.)
4) $\vec{MA}+2\vec{MB}+\vec{MC}=4\vec{MG}$, donc $\|4\vec{MG}\|=8 \iff \|\vec{MG}\|=2$ :
$(\mathcal{E})$ est le cercle de centre $G\big(\tfrac72\,;\,1\big)$ et de rayon~$2$.

\smallskip
\noindent\textbf{Partie B — Situation.}
\emph{Tâche 1.} $C(x)=x^2-8x+25=(x-4)^2+9$. Le coût est minimal pour $x=4$ (centaines de
livres, soit $400$ livres) et vaut $C(4)=9$ milliers de francs.

\emph{Tâche 2.} $u_n=500\times1{,}15^{\,n}$. On résout $u_n>1000 \iff 1{,}15^{\,n}>2
\iff n>\dfrac{\ln 2}{\ln 1{,}15}\approx4{,}96$ : le nombre d'abonnés dépasse $1000$ à partir
de la \textbf{5\textsuperscript{e} année} ($u_5=500\times1{,}15^5\approx1006$).

\emph{Tâche 3.} Le nombre de livres mal collés suit une loi binomiale $\mathcal{B}(3\,;\,0{,}04)$.
$P(\text{au moins un})=1-(0{,}96)^3=1-0{,}884736\approx\mathbf{0{,}115}$, soit environ
\SI{11,5}{\percent}.
